%2multibyte Version: 5.50.0.2960 CodePage: 65001
% !TeX root = ./Slides_Lez_16_Python_Programmazione_Stocastica.tex
% ------------------------------------------------------------
% Macro matematiche ufficiali del progetto MQF
% Coerenti con MQF_Notazione.tex
% ------------------------------------------------------------
% Fallback utile per compilazione esterna se qualche titolo resta scritto come \QTR.
% Scientific WorkPlace usa \QTR internamente; in Beamer \QTR{frametitle}{Titolo}
% deve espandersi in \frametitle{Titolo}.

\documentclass[handout]{beamer}%
\usepackage{mathpazo}
%\usepackage{hyperref}
\usepackage{multimedia}
\usepackage{amsmath}
\usepackage{amsfonts}
\usepackage{amssymb}
\usepackage{graphicx}%
\setcounter{MaxMatrixCols}{30}
%TCIDATA{OutputFilter=latex2.dll}
%TCIDATA{Version=5.50.0.2960}
%TCIDATA{Codepage=65001}
%TCIDATA{CSTFile=beamer.cst}
%TCIDATA{Created=Saturday, May 16, 2026 09:09:59}
%TCIDATA{LastRevised=Saturday, May 16, 2026 11:24:58}
%TCIDATA{<META NAME="GraphicsSave" CONTENT="32">}
%TCIDATA{<META NAME="SaveForMode" CONTENT="2">}
%TCIDATA{BibliographyScheme=Manual}
%TCIDATA{<META NAME="DocumentShell" CONTENT="Other Documents\SW\Slides - Beamer">}
%BeginMSIPreambleData
\graphicspath{{../graphics/}}
\providecommand{\U}[1]{\protect\rule{.1in}{.1in}}
%EndMSIPreambleData
\newenvironment{stepenumerate}{\begin{enumerate}[<+->]}{\end{enumerate}}
\newenvironment{stepitemize}{\begin{itemize}[<+->]}{\end{itemize} }
\newenvironment{stepenumeratewithalert}{\begin{enumerate}[<+-| alert@+>]}{\end{enumerate}}
\newenvironment{stepitemizewithalert}{\begin{itemize}[<+-| alert@+>]}{\end{itemize} }
\usetheme{Madrid}


\setbeamertemplate{footline}
{
  \leavevmode%
  \hbox{%
    \begin{beamercolorbox}[wd=.20\paperwidth,ht=2.25ex,dp=1ex,center]
      {author in head/foot}%
      \usebeamerfont{author in head/foot}\insertshortauthor
    \end{beamercolorbox}%
    \begin{beamercolorbox}[wd=.60\paperwidth,ht=2.25ex,dp=1ex,center]
      {title in head/foot}%
      \usebeamerfont{title in head/foot}\insertshorttitle
    \end{beamercolorbox}%
    \begin{beamercolorbox}[wd=.20\paperwidth,ht=2.25ex,dp=1ex,right]
      {date in head/foot}%
      \usebeamerfont{date in head/foot}%
      \insertshortdate{}\hspace*{2em}%
      \insertframenumber{} / \inserttotalframenumber\hspace*{2ex}
    \end{beamercolorbox}%
  }%
  \vskip0pt%
}




\providecommand{\R}{\mathbb{R}}
\providecommand{\N}{\mathbb{N}}
\providecommand{\E}{\mathbb{E}}
\providecommand{\Prob}{\mathbb{P}}
\providecommand{\Q}{\mathbb{Q}}
\providecommand{\F}{\mathcal{F}}
\providecommand{\B}{\mathcal{B}}
\providecommand{\G}{\mathcal{G}}
\providecommand{\1}{\mathbf{1}}
\providecommand{\Var}{\operatorname{Var}}
\providecommand{\Cov}{\operatorname{Cov}}
\providecommand{\VaR}{\operatorname{VaR}}
\providecommand{\CVaR}{\operatorname{CVaR}}
\providecommand{\argmin}{\operatorname*{arg\,min}}
\providecommand{\argmax}{\operatorname*{arg\,max}}
\providecommand{\EVPI}{\operatorname{EVPI}}
\providecommand{\VSS}{\operatorname{VSS}}
\providecommand{\QTR}[2]{\csname #1\endcsname{#2}}
%BeginMSIPreambleData
\ifx\pdfoutput\relax\let\pdfoutput=\undefined\fi
\newcount\msipdfoutput
\ifx\pdfoutput\undefined\else
\ifcase\pdfoutput\else
\msipdfoutput=1
\ifx\paperwidth\undefined\else
\ifdim\paperheight=0pt\relax\else\pdfpageheight\paperheight\fi
\ifdim\paperwidth=0pt\relax\else\pdfpagewidth\paperwidth\fi
\fi\fi\fi
%EndMSIPreambleData

%******************
% Slide 1
%******************

\title[Lez. 16 - Python e programmazione stocastica]
{Lezione 16 - Applicazione in Python: programmazione stocastica}

\subtitle{Metodi Quantitativi per la Finanza}

\author[P. Falbo, C. Pelizzari]
{{\LARGE Paolo Falbo, Cristian Pelizzari}}

\institute[]
{{\large Dipartimento di Economia e Management}}

\date[]{}

\titlegraphic{
\includegraphics[width=3.3cm]{../graphics/logo-Unibs-esteso.png}
}

\begin{document}

\maketitle


\section{Apertura}

%******************
% Slide 2
%******************

\subsection{Obiettivi e struttura della lezione}

%TCIMACRO{\TeXButton{BeginFrame}{\begin{frame}}}%
%BeginExpansion
\begin{frame}%
%EndExpansion

\QTR{frametitle}{Obiettivi e struttura della lezione}

%TCIMACRO{\TeXButton{Transparency}{\setbeamercovered{transparent=20}}}%
%BeginExpansion
\setbeamercovered{transparent=20}%
%EndExpansion

La Lezione 16 conclude il percorso delle lezioni applicative del corso.

%\vspace{0.3cm}

Gli obiettivi sono tre.

\begin{stepenumerate}

\item [1.] \textbf{Focus Python}

Richiamare come un problema di ottimizzazione lineare viene tradotto nella
struttura numerica richiesta da \texttt{scipy.optimize.linprog}.

\item [2.] \textbf{Caso Aula: SVB--ALM 2023}

Tradurre i modelli dei Capitoli 14 e 15 in un problema computazionale
two-stage e confrontare soluzione stocastica, Expected Value,
Expected result of using the EV solution e Wait-and-See.

\item [3.] \textbf{Caso TakeHome: UK LDI 2022}

Trasferire autonomamente lo stesso metodo a un problema di buffer di
liquidit\`a, margin call e deleveraging.

\end{stepenumerate}

%\vspace{0.3cm}

Il metodo resta quello introdotto nella Lezione 4:
%\vspace{-0.1cm}
\[
\text{modello}
\;\longrightarrow\;
\text{implementazione}
\;\longrightarrow\;
\text{controllo}
\;\longrightarrow\;
\text{interpretazione}.
\]

%\vspace{0.2cm}

Dopo il break, il Caso Aula sar\`a sviluppato direttamente in
\textbf{VS Code/Jupyter}, senza ulteriori slide.

%TCIMACRO{\TeXButton{Transition: Box Out}{\transboxout}}%
%BeginExpansion
\transboxout%
%EndExpansion
%TCIMACRO{\TeXButton{EndFrame}{\end{frame}}}%
%BeginExpansion
\end{frame}%
%EndExpansion


\section{Focus Python}

%******************
% Slide 3
%******************

\subsection{Dal modello matematico al solver}

%TCIMACRO{\TeXButton{BeginFrame}{\begin{frame}}}%
%BeginExpansion
\begin{frame}%
%EndExpansion

\QTR{frametitle}{Dal modello matematico a \texttt{linprog}}

%TCIMACRO{\TeXButton{Transparency}{\setbeamercovered{transparent=20}}}%
%BeginExpansion
\setbeamercovered{transparent=20}%
%EndExpansion

Un solver lineare non ``vede'' il significato finanziario delle variabili.

%\vspace{0.3cm}

Riceve una rappresentazione numerica del problema:

\[
\begin{array}{ll}
\max\limits_x & c'x \\[0.2cm]
\text{s.t.}   & A_{\mathrm{eq}}x=b_{\mathrm{eq}},\\
              & A_{\mathrm{ub}}x\le b_{\mathrm{ub}},\\
              & \ell\le x\le u.
\end{array}
\]

%\vspace{0.3cm}

In \texttt{scipy.optimize.linprog}, la corrispondenza \`e:

\[
\begin{array}{ccl}
c & \longrightarrow & \texttt{c},\\
A_{\mathrm{eq}},\,b_{\mathrm{eq}}
& \longrightarrow &
\texttt{A\_eq},\,\texttt{b\_eq},\\
A_{\mathrm{ub}},\,b_{\mathrm{ub}}
& \longrightarrow &
\texttt{A\_ub},\,\texttt{b\_ub},\\
\ell,\,u
& \longrightarrow &
\texttt{bounds}.
\end{array}
\]

%\vspace{0.3cm}

\textbf{Punto chiave:} costruire vettori e matrici significa tradurre il
modello matematico in una forma computazionale coerente.

%TCIMACRO{\TeXButton{Transition: Box Out}{\transboxout}}%
%BeginExpansion
\transboxout%
%EndExpansion
%TCIMACRO{\TeXButton{EndFrame}{\end{frame}}}%
%BeginExpansion
\end{frame}%
%EndExpansion

%******************
% Slide 4
%******************

\subsection{Ordinamento delle variabili e funzione obiettivo}

%TCIMACRO{\TeXButton{BeginFrame}{\begin{frame}}}%
%BeginExpansion
\begin{frame}%
%EndExpansion

\QTR{frametitle}{Due dettagli che il solver non pu\`o correggere}

%TCIMACRO{\TeXButton{Transparency}{\setbeamercovered{transparent=20}}}%
%BeginExpansion
\setbeamercovered{transparent=20}%
%EndExpansion

Nel Caso Aula, useremo un vettore decisionale ordinato come

\[
z=
(x_1,x_2,x_3,x_4,
u_N,u_P,u_R,
g_N,g_P,g_R)'.
\]

\vspace{0.3cm}

\begin{stepenumerate}

\item [1.] \textbf{L'ordine delle variabili conta}

La colonna \(j\) delle matrici deve riferirsi alla componente
\(z_j\) del vettore.

Lo stesso ordine riguarda gli output del solver.

\vspace{0.2cm}

\item [2.] \textbf{\texttt{linprog} risolve un problema di minimizzazione}

Se il modello finanziario richiede $\max_z f(z)$, al solver viene passato l'obiettivo equivalente
%\vspace{-0.2cm}
\[
\min_z \{-f(z)\}.
\]

\end{stepenumerate}

\vspace{0.3cm}

Nel codice, quindi, il segno dei coefficienti dell'obiettivo non \`e una
convenzione estetica: deve essere coerente con il problema originario.

%\vspace{0.2cm}

Un primo controllo tecnico sar\`a sempre: \texttt{res.success}

%TCIMACRO{\TeXButton{Transition: Box Out}{\transboxout}}%
%BeginExpansion
\transboxout%
%EndExpansion
%TCIMACRO{\TeXButton{EndFrame}{\end{frame}}}%
%BeginExpansion
\end{frame}%
%EndExpansion


\section{Richiamo dei Capitoli 14--15}

%******************
% Slide 5
%******************

\subsection{La struttura two-stage}

%TCIMACRO{\TeXButton{BeginFrame}{\begin{frame}}}%
%BeginExpansion
\begin{frame}%
%EndExpansion

\QTR{frametitle}{Prima si decide, poi si osserva, infine si reagisce}

%TCIMACRO{\TeXButton{Transparency}{\setbeamercovered{transparent=20}}}%
%BeginExpansion
\setbeamercovered{transparent=20}%
%EndExpansion

Nel Capitolo 14, abbiamo introdotto la struttura fondamentale di un problema
stocastico a due stadi:
\vspace{-0.1cm}
\[
x
\;\longrightarrow\;
s
\;\longrightarrow\;
y_s.
\]
\vspace{-0.5cm}
%\vspace{0.3cm}

\begin{stepitemize}

\item \(x\) \`e la \textbf{decisione di primo stadio}, assunta prima di
conoscere lo scenario futuro.

\item \(s\in\mathcal S\) \`e lo \textbf{scenario} che si realizza dopo la
decisione iniziale.

\item \(y_s\) \`e la \textbf{decisione di ricorso}, scelta dopo avere
osservato lo scenario \(s\).

\end{stepitemize}

\vspace{0.3cm}

La funzione obiettivo combina quindi il risultato della decisione iniziale e
il valore atteso delle decisioni di ricorso:

\[
\max_x
\left\{
c'x+\sum_{s\in\mathcal S}p_s Q(x,s)
\right\}.
\]

\vspace{0.2cm}

\textbf{Idea centrale:} la decisione iniziale deve essere robusta rispetto
agli scenari; il ricorso utilizza l'informazione disponibile dopo la loro
realizzazione.

%TCIMACRO{\TeXButton{Transition: Box Out}{\transboxout}}%
%BeginExpansion
\transboxout%
%EndExpansion
%TCIMACRO{\TeXButton{EndFrame}{\end{frame}}}%
%BeginExpansion
\end{frame}%
%EndExpansion



%******************
% Slide 6
%******************

\subsection{Informazione e non anticipativit\`a}

%TCIMACRO{\TeXButton{BeginFrame}{\begin{frame}}}%
%BeginExpansion
\begin{frame}%
%EndExpansion

\QTR{frametitle}{La decisione iniziale non conosce ancora lo scenario}

%TCIMACRO{\TeXButton{Transparency}{\setbeamercovered{transparent=20}}}%
%BeginExpansion
\setbeamercovered{transparent=20}%
%EndExpansion

Nel problema two-stage, la decisione iniziale \(x\) viene presa
\textbf{prima} di conoscere quale scenario \(s\in\mathcal S\) si realizzer\`a.

%\vspace{0.4cm}

Per questo motivo \(x\), non pu\`o dipendere dallo scenario futuro:
\vspace{-0.0cm}
\[
\boxed{
x
\quad\text{\`e unica e viene scelta ex ante}
}
\]

%\vspace{0.4cm}

Dopo la realizzazione dello scenario, invece, la decisione di ricorso pu\`o
adattarsi all'informazione osservata:
\vspace{-0.1cm}
\[
y_N,\qquad y_P,\qquad y_R.
\]

%\vspace{0.4cm}

La struttura informativa \`e quindi
\vspace{-0.1cm}
\[
\underbrace{x}_{\text{prima di conoscere }s}
\quad
\longrightarrow
\quad
\underbrace{s}_{\text{scenario osservato}}
\quad
\longrightarrow
\quad
\underbrace{y_s}_{\text{ricorso adattato a }s}.
\]

%\vspace{0.3cm}

\textbf{Non anticipativit\`a:} la decisione di primo stadio non pu\`o
utilizzare informazione che sar\`a disponibile soltanto dopo.

%TCIMACRO{\TeXButton{Transition: Box Out}{\transboxout}}%
%BeginExpansion
\transboxout%
%EndExpansion
%TCIMACRO{\TeXButton{EndFrame}{\end{frame}}}%
%BeginExpansion
\end{frame}%
%EndExpansion


%******************
% Slide 7
%******************

\subsection{Benchmark e valore dell'informazione}

%TCIMACRO{\TeXButton{BeginFrame}{\begin{frame}}}%
%BeginExpansion
\begin{frame}%
%EndExpansion

\QTR{frametitle}{Tre benchmark, due misure di valore}

%TCIMACRO{\TeXButton{Transparency}{\setbeamercovered{transparent=20}}}%
%BeginExpansion
\setbeamercovered{transparent=20}%
%EndExpansion

Il Capitolo 15 confronta tre modi diversi di utilizzare l'informazione.

%\vspace{0.2cm}

\begin{description}

\item[\textbf{SP}]
si sceglie oggi utilizzando l'intera distribuzione degli scenari.

\item[\textbf{EV/EEV}]
si risolve prima il problema costruito sui valori medi; la decisione
\(x^{EV}\) viene poi fissata e rivalutata negli scenari originari.

\item[\textbf{WS}]
si assume, come benchmark, di conoscere lo scenario prima di scegliere.

\end{description}

%\vspace{0.3cm}

Per un problema di massimizzazione, quando le politiche confrontate sono
ammissibili,
\vspace{-0.0cm}
\[
z^{EEV}
\le
z^{SP}
\le
z^{WS}.
\]

%\vspace{0.2cm}

Le due misure principali sono
\vspace{-0.0cm}
\[
VSS=z^{SP}-z^{EEV},
\qquad
EVPI=z^{WS}-z^{SP}.
\]

%\vspace{0.3cm}

\textbf{Attenzione:} il valore ottimo del problema costruito sui parametri
medi non \`e \(z^{EEV}\). Per ottenere \(z^{EEV}\), \(x^{EV}\) deve essere
rivalutata nella distribuzione originaria degli scenari.

%TCIMACRO{\TeXButton{Transition: Box Out}{\transboxout}}%
%BeginExpansion
\transboxout%
%EndExpansion
%TCIMACRO{\TeXButton{EndFrame}{\end{frame}}}%
%BeginExpansion
\end{frame}%
%EndExpansion



\section{Caso Aula: modello}

%******************
% Slide 8
%******************

\subsection{Problema finanziario e decisione iniziale}

%TCIMACRO{\TeXButton{BeginFrame}{\begin{frame}}}%
%BeginExpansion
\begin{frame}%
%EndExpansion

\QTR{frametitle}{SVB--ALM 2023: il problema}

%TCIMACRO{\TeXButton{Transparency}{\setbeamercovered{transparent=20}}}%
%BeginExpansion
\setbeamercovered{transparent=20}%
%EndExpansion

Consideriamo una banca che dispone inizialmente di
\vspace{-0.0cm}
\[
A_0=100
\]

unit\`a monetarie e deve scegliere la composizione dell'attivo
\textbf{prima} di conoscere lo scenario futuro di liquidit\`a.

%\vspace{0.4cm}

La domanda quantitativa \`e:

\begin{center}
\textbf{come allocare le 100 unit\`a per massimizzare il risultato economico
atteso, tenendo conto del rischio di fabbisogno di liquidit\`a?}
\end{center}

%\vspace{0.4cm}

Il problema combina tre elementi:

\begin{stepitemize}

\item rendimento delle attivit\`a;

\item capacit\`a delle attivit\`a di generare liquidit\`a nei diversi scenari;

\item costo del funding di emergenza.

\end{stepitemize}

\vspace{0.3cm}

Il caso \`e un modello stilizzato ispirato all'episodio SVB del 2023.

%TCIMACRO{\TeXButton{Transition: Box Out}{\transboxout}}%
%BeginExpansion
\transboxout%
%EndExpansion
%TCIMACRO{\TeXButton{EndFrame}{\end{frame}}}%
%BeginExpansion
\end{frame}%
%EndExpansion


%******************
% Slide 9
%******************

\subsection{Le decisioni di primo stadio}

%TCIMACRO{\TeXButton{BeginFrame}{\begin{frame}}}%
%BeginExpansion
\begin{frame}%
%EndExpansion

\QTR{frametitle}{Quattro classi di attivo}

%TCIMACRO{\TeXButton{Transparency}{\setbeamercovered{transparent=20}}}%
%BeginExpansion
\setbeamercovered{transparent=20}%
%EndExpansion

La decisione iniziale \`e
\vspace{-0.1cm}
\[
x=(x_1,x_2,x_3,x_4)'.
\]

\vspace{-0.2cm}

\[
\begin{array}{c|l|c}
\text{Variabile} & \text{Classe di attivo} & c_j \\ \hline
x_1 & \text{cassa e riserve} & 0.010 \\
x_2 & \text{titoli a breve} & 0.025 \\
x_3 & \text{titoli a lunga} & 0.040 \\
x_4 & \text{prestiti e impieghi meno liquidi} & 0.050
\end{array}
\]

\vspace{+0.2cm}

Il vincolo di bilancio iniziale \`e
\vspace{-0.0cm}
\[
x_1+x_2+x_3+x_4=100,
\]
\vspace{-0.0cm}
con
\vspace{-0.0cm}
\[
x_j\ge0,
\qquad
x_3\le70,
\qquad
x_4\le70.
\]

%\vspace{0.3cm}

Le quattro classi di attivo offrono rendimenti diversi. In seguito, vedremo che differiscono anche per capacit\`a di generare
liquidit\`a nei diversi scenari: la scelta iniziale dovr\`a quindi bilanciare
rendimento e vulnerabilit\`a al fabbisogno di liquidit\`a.

%TCIMACRO{\TeXButton{Transition: Box Out}{\transboxout}}%
%BeginExpansion
\transboxout%
%EndExpansion
%TCIMACRO{\TeXButton{EndFrame}{\end{frame}}}%
%BeginExpansion
\end{frame}%
%EndExpansion


%******************
% Slide 10
%******************

\subsection{Gli scenari di liquidit\`a}

%TCIMACRO{\TeXButton{BeginFrame}{\begin{frame}}}%
%BeginExpansion
\begin{frame}%
%EndExpansion

\QTR{frametitle}{Tre scenari, tre fabbisogni}

%TCIMACRO{\TeXButton{Transparency}{\setbeamercovered{transparent=20}}}%
%BeginExpansion
\setbeamercovered{transparent=20}%
%EndExpansion

Dopo la decisione iniziale pu\`o realizzarsi uno dei tre scenari

\[
\mathcal S=\{N,P,R\}.
\]

\vspace{0.3cm}

\[
\begin{array}{c|l|c|c}
s & \text{Scenario} & p_s & d_s \\ \hline
N & \text{condizioni normali} & 0.60 & 15 \\
P & \text{pressione sulla raccolta} & 0.25 & 35 \\
R & \text{run severo} & 0.15 & 65
\end{array}
\]

\vspace{0.4cm}

\(d_s\) misura il fabbisogno di liquidit\`a nello scenario \(s\).

\vspace{0.3cm}

L'incertezza riguarda quindi non soltanto \textbf{quanta} liquidit\`a serve,
ma anche \textbf{quanta} liquidit\`a gli attivi riescono a generare nello
stesso scenario.

%TCIMACRO{\TeXButton{Transition: Box Out}{\transboxout}}%
%BeginExpansion
\transboxout%
%EndExpansion
%TCIMACRO{\TeXButton{EndFrame}{\end{frame}}}%
%BeginExpansion
\end{frame}%
%EndExpansion


%******************
% Slide 11
%******************

\subsection{Liquidabilit\`a scenario-specifica}

%TCIMACRO{\TeXButton{BeginFrame}{\begin{frame}}}%
%BeginExpansion
\begin{frame}%
%EndExpansion

\QTR{frametitle}{Rendimento e liquidabilit\`a sono grandezze diverse}

%TCIMACRO{\TeXButton{Transparency}{\setbeamercovered{transparent=20}}}%
%BeginExpansion
\setbeamercovered{transparent=20}%
%EndExpansion

Indichiamo con
\vspace{-0.1cm}
\[
\alpha_{js}
\]

la quota dell'attivo \(j\) trasformabile in liquidit\`a nello scenario \(s\).

%\vspace{0.3cm}

La liquidit\`a ottenibile dal portafoglio nello scenario \(s\) \`e
\vspace{-0.0cm}
\[
L_s(x)
=
\sum_{j=1}^{4}\alpha_{js}x_j.
\]

%\vspace{0.4cm}

Per esempio, nello scenario \(R\):
\vspace{-0.0cm}
\[
(\alpha_{1R},\alpha_{2R},\alpha_{3R},\alpha_{4R})
=
(1.00,0.90,0.50,0.20).
\]

%\vspace{0.4cm}

Occorre distinguere nettamente:

\[
\underbrace{c_j}_{\text{rendimento unitario}}
\qquad\neq\qquad
\underbrace{\alpha_{js}}_{\text{quota liquidabile nello scenario }s}.
\]

%\vspace{0.3cm}

Un'attivit\`a pu\`o essere redditizia ex ante, ma poco efficace nel generare
liquidit\`a durante uno stress.

%TCIMACRO{\TeXButton{Transition: Box Out}{\transboxout}}%
%BeginExpansion
\transboxout%
%EndExpansion
%TCIMACRO{\TeXButton{EndFrame}{\end{frame}}}%
%BeginExpansion
\end{frame}%
%EndExpansion


%******************
% Slide 12
%******************

\subsection{Il ricorso}

%TCIMACRO{\TeXButton{BeginFrame}{\begin{frame}}}%
%BeginExpansion
\begin{frame}%
%EndExpansion

\QTR{frametitle}{Dopo lo scenario: funding e liquidit\`a residua}

%TCIMACRO{\TeXButton{Transparency}{\setbeamercovered{transparent=20}}}%
%BeginExpansion
\setbeamercovered{transparent=20}%
%EndExpansion

Dopo avere osservato lo scenario \(s\), la banca pu\`o utilizzare due
variabili di ricorso:
\vspace{-0.0cm}
\[
u_s\ge0
\qquad
\text{funding di emergenza},
\]
\vspace{-0.0cm}
\[
g_s\ge0
\qquad
\text{liquidit\`a residua}.
\]

%\vspace{0.4cm}

Il bilancio di liquidit\`a scenario-specifico \`e
\vspace{-0.0cm}
\[
\sum_{j=1}^{4}\alpha_{js}x_j+u_s
=
d_s+g_s.
\]

%\vspace{0.4cm}

Tutti i termini sono espressi nella stessa unit\`a monetaria:
\vspace{-0.0cm}
\[
\underbrace{\sum_j\alpha_{js}x_j}_{\text{liquidit\`a dagli attivi}}
+
\underbrace{u_s}_{\text{funding}}
=
\underbrace{d_s}_{\text{fabbisogno}}
+
\underbrace{g_s}_{\text{residuo}}.
\]

\vspace{0.3cm}

Il ricorso consente quindi di adattare la risposta della banca allo scenario
osservato, mantenendo invariata la decisione iniziale \(x\).

%TCIMACRO{\TeXButton{Transition: Box Out}{\transboxout}}%
%BeginExpansion
\transboxout%
%EndExpansion
%TCIMACRO{\TeXButton{EndFrame}{\end{frame}}}%
%BeginExpansion
\end{frame}%
%EndExpansion

```latex
%******************
% Slide 13
%******************

\subsection{La funzione obiettivo}

%TCIMACRO{\TeXButton{BeginFrame}{\begin{frame}}}%
%BeginExpansion
\begin{frame}%
%EndExpansion

\QTR{frametitle}{Rendimento contro costo atteso del ricorso}

%TCIMACRO{\TeXButton{Transparency}{\setbeamercovered{transparent=20}}}%
%BeginExpansion
\setbeamercovered{transparent=20}%
%EndExpansion

Il rendimento iniziale del portafoglio \`e

\[
\sum_{j=1}^{4}c_jx_j.
\]

\vspace{0.3cm}

Il funding di emergenza nello scenario \(s\) ha costo unitario
\(\kappa_s\), quindi il costo scenario-specifico \`e

\[
\kappa_su_s.
\]

\vspace{0.3cm}

Il programma stocastico massimizza

\[
\boxed{
\sum_{j=1}^{4}c_jx_j
-
\sum_{s\in\mathcal S}p_s\kappa_su_s
}
\]

soggetto ai vincoli di allocazione e ai bilanci di liquidit\`a.

\vspace{0.4cm}

La funzione obiettivo separa due componenti:

\[
\underbrace{\sum_j c_jx_j}_{\text{rendimento iniziale}}
-
\underbrace{\sum_s p_s\kappa_su_s}_{\text{costo atteso del funding}}.
\]

\vspace{0.3cm}

\textbf{Interpretazione:} il modello ricerca rendimento, ma attribuisce un
costo economico alla dipendenza dal funding di emergenza negli scenari pi\`u
critici.

%TCIMACRO{\TeXButton{Transition: Box Out}{\transboxout}}%
%BeginExpansion
\transboxout%
%EndExpansion
%TCIMACRO{\TeXButton{EndFrame}{\end{frame}}}%
%BeginExpansion
\end{frame}%
%EndExpansion
```








\end{document}